\section{Phase I: Dataset Freeze and Version-Window Definition}
\label{sec:phase1_dataset_curation}

This section documents the completed Phase I freeze workflow for the centralized dataset:
repository metadata freeze, stable-tag freeze, and consecutive-window construction. The first
four repositories were used as a pilot to validate the workflow. The final Phase I dataset
keeps those four repositories and adds six more projects from different ecosystems. The cohort
labels are retained only as provenance; the analysis dataset is centralized.

\subsection{Repository Freeze}
\label{sec:phase1_frozen_repos}

The centralized Phase I dataset contains ten mature GitHub-hosted repositories. Storage
persistence is recorded as \texttt{High (GitHub)} for every project. Software-type tags are
kept short and descriptive so that later analysis can group projects without losing the
identity of each repository.

\begin{table}[H]
\centering
\scriptsize
\renewcommand{\arraystretch}{1.12}
\begin{tabularx}{\textwidth}{p{1.25cm}p{1.45cm}X p{3.2cm}}
\toprule
\textbf{Cohort} & \textbf{Project} & \textbf{Repository URL} & \textbf{Software-type tags} \\
\midrule
Pilot & flask    & \url{https://github.com/pallets/flask}     & web-framework, library \\
Pilot & requests & \url{https://github.com/psf/requests}      & http-client, library \\
Pilot & urllib3  & \url{https://github.com/urllib3/urllib3}   & http-client, library \\
Pilot & click    & \url{https://github.com/pallets/click}     & cli-toolkit, library \\
Extended & brew     & \url{https://github.com/Homebrew/brew}     & package-manager, cli-tool, ecosystem-infrastructure \\
Extended & curl     & \url{https://github.com/curl/curl}         & network-transfer-tool, cli-tool, library \\
Extended & gin      & \url{https://github.com/gin-gonic/gin}     & web-framework, http-framework, library \\
Extended & junit5   & \url{https://github.com/junit-team/junit5} & testing-framework, library \\
Extended & prettier & \url{https://github.com/prettier/prettier} & code-formatter, cli-tool, developer-tool \\
Extended & tokio    & \url{https://github.com/tokio-rs/tokio}    & async-runtime, concurrency-library, library \\
\bottomrule
\end{tabularx}
\caption{Centralized Phase I repository freeze.}
\label{tab:round1_freeze}
\end{table}

\subsection{Tag Freeze and Consecutive Windows}
\label{sec:phase1_frozen_windows}

For each project, eight stable release tags are frozen. Let $\{T_1,\dots,T_8\}$ be the
frozen tags sorted chronologically by tag date from oldest to newest. Seven tag-chronological windows are then defined as
$W_k=(T_k \rightarrow T_{k+1})$ for $k\in\{1,\dots,7\}$. Each frozen tag is bound to a
specific commit hash, including dereferencing annotated tags to their target commits. This
gives each metric value a replayable repository state and deterministic Git revision expression. It does not by itself
guarantee that every adjacent tag pair belongs to one ancestor-to-descendant release line; that condition is evaluated
separately in Section~\ref{sec:posthoc_ancestry_validation}.

The final freeze is captured in \path{data/dataset_freeze_all.csv}, \path{data/versions_all.csv},
and \path{data/windows_all.csv}. The earlier pilot-only CSVs were consolidated into these centralized
files and removed from the active evidence layer to avoid maintaining parallel copies of the same dataset.

Centralized Phase I totals are:
\begin{itemize}[leftmargin=*]
  \item 10 repositories;
  \item 80 frozen tags (10 projects $\times$ 8 tags);
  \item 70 consecutive windows (10 projects $\times$ 7 windows);
  \item tag-date coverage from 2022-06-06 to 2026-04-29.
\end{itemize}

\begin{table}[H]
\centering
\scriptsize
\renewcommand{\arraystretch}{1.12}
\begin{tabular}{lrrll}
\toprule
\textbf{Project} & \textbf{Tags} & \textbf{Windows} & \textbf{First tag} & \textbf{Last tag} \\
\midrule
brew     & 8 & 7 & 5.1.1        & 5.1.8 \\
click    & 8 & 7 & 8.1.6        & 8.3.1 \\
curl     & 8 & 7 & curl-8\_14\_0  & curl-8\_20\_0 \\
flask    & 8 & 7 & 3.0.0        & 3.1.3 \\
gin      & 8 & 7 & v1.8.1       & v1.12.0 \\
junit5   & 8 & 7 & r6.0.0       & r5.14.4 \\
prettier & 8 & 7 & 3.7.1        & 3.8.3 \\
requests & 8 & 7 & v2.30.0      & v2.32.5 \\
tokio    & 8 & 7 & tokio-1.47.3 & tokio-1.52.1 \\
urllib3  & 8 & 7 & 2.2.3        & 2.6.3 \\
\bottomrule
\end{tabular}
\caption{Centralized version coverage by project.}
\label{tab:round1_version_coverage}
\end{table}

\begin{table}[H]
\centering
\scriptsize
\renewcommand{\arraystretch}{1.12}
\begin{tabularx}{\textwidth}{p{1.7cm}p{3.6cm}X}
\toprule
\textbf{Project} & \textbf{Stable tag rule} & \textbf{Reason for project-specific rule} \\
\midrule
brew & \path{^\d+\.\d+\.\d+$} & Plain semantic-version release tags. \\
click & \path{^\d+\.\d+\.\d+$} & Plain semantic-version release tags. \\
curl & \path{^curl-\d+_\d+_\d+$} & Stable curl releases use a \texttt{curl-} prefix and underscores; release candidates are excluded. \\
flask & \path{^\d+\.\d+\.\d+$} & Plain semantic-version release tags. \\
gin & \path{^v\d+\.\d+\.\d+$} & Stable Go module tags use a \texttt{v} prefix. \\
junit5 & \path{^r\d+\.\d+\.\d+$} & Stable JUnit tags use an \texttt{r} prefix; milestone and release-candidate tags are excluded. \\
prettier & \path{^\d+\.\d+\.\d+$} & Stable releases use plain semantic versions; alpha and non-version tags are excluded. \\
requests & \path{^v\d+\.\d+\.\d+$} & Stable release tags use a \texttt{v} prefix. \\
tokio & \path{^tokio-\d+\.\d+\.\d+$} & The repository contains crate-specific tags; only main Tokio runtime releases are included. \\
urllib3 & \path{^\d+\.\d+\.\d+$} & Plain semantic-version release tags. \\
\bottomrule
\end{tabularx}
\caption{Stable-tag rules for the centralized repositories.}
\label{tab:centralized_tag_rules}
\end{table}

\begin{table}[H]
\centering
\scriptsize
\renewcommand{\arraystretch}{1.12}
\begin{tabular}{lrrr}
\toprule
\textbf{Project} & \textbf{Min window days} & \textbf{Max window days} & \textbf{Mean window days} \\
\midrule
brew     & 0.83 & 8.96   & 4.99 \\
click    & 10.04 & 489.99 & 121.58 \\
curl     & 7.00 & 63.00  & 48.00 \\
flask    & 16.05 & 219.96 & 124.66 \\
gin      & 60.82 & 378.26 & 194.70 \\
junit5   & 0.03 & 69.91  & 29.73 \\
prettier & 0.87 & 78.36  & 19.79 \\
requests & 0.26 & 376.04 & 119.74 \\
tokio    & 0.58 & 58.99  & 14.87 \\
urllib3  & 3.00 & 170.05 & 68.89 \\
\bottomrule
\end{tabular}
\caption[Window-duration statistics.]{Window-duration statistics for the centralized release windows.}
\label{tab:round1_window_durations}
\end{table}

Two details need to be kept in view during Phase II. First, some repositories release several
stable tags very close together, especially \texttt{junit5}, \texttt{brew}, \texttt{prettier},
and \texttt{tokio}. Second, JUnit has overlapping stable lines in the frozen interval. This is
preserved in the original tag-date freeze for reproducibility, but it cannot be handled by short-window flags alone.
The post-hoc ancestry audit therefore identifies divergent pairs explicitly, excludes them from the ancestry-only C1
sensitivity analysis, and makes them ineligible for the baseline PoC.
